\documentclass[11pt]{article}
\usepackage[a4paper,margin=2.6cm]{geometry}
\usepackage{amsmath,amssymb,amsthm,mathtools}
\usepackage{booktabs}
\usepackage{enumitem}
\usepackage[hidelinks]{hyperref}

\newtheorem{theorem}{Theorem}[section]
\newtheorem{lemma}[theorem]{Lemma}
\newtheorem{proposition}[theorem]{Proposition}
\newtheorem{corollary}[theorem]{Corollary}
\theoremstyle{definition}
\newtheorem{definition}[theorem]{Definition}
\newtheorem{remark}[theorem]{Remark}
\newtheorem{protocol}[theorem]{Protocol}

\newcommand{\F}{\mathbb{F}}
\newcommand{\K}{K}
\newcommand{\RS}{\mathrm{RS}}
\newcommand{\Fold}{\mathrm{Fold}}
\newcommand{\rev}{\mathrm{rev}}
\newcommand{\wt}{\mathrm{wt}}
\newcommand{\Bad}{\mathsf{Bad}}
\newcommand{\Win}{\mathsf{Win}}
\newcommand{\Acc}{\mathsf{Acc}}
\DeclareMathOperator{\dist}{d}

\title{Hypercube Sums via a Dual Kronecker Chain:\\
Sumcheck on Growing Domains, Polynomial-Valued FRI on Shrinking Domains}
\author{Abdelali Mkhida\\ \small Algorizk}
\date{November 2026}

\begin{document}
\maketitle

\begin{abstract}
We study an interactive oracle proof for the Boolean-hypercube sum $H=\sum_{b\in\{0,1\}^m} f(b)$ of a committed multilinear polynomial $f$. The protocol runs two synchronized chains on the same challenges $r_1,\dots,r_m$: a forward Kronecker--Sumcheck chain on growing multiplicative domains $D_0\subset\cdots\subset D_m$, and a backward polynomial-valued (symbolic) FRI chain on shrinking domains, whose leaves are polynomials in a formal variable $X$. Both chains terminate at $f(r_1X^{2^{m-1}},\dots,r_mX)$. We prove completeness for every $m$, show that in the exact-identity model both chains are deterministic bijections, give two explicit attacks on under-specified variants (one accepting a false sum with probability $1$), and prove a soundness theorem for the repaired protocol with the explicit bound $m/|\K|+2B(2^m-1)/|\K|+(1-\theta)^q$, $\theta=(1-1/B)/2$. For Goldilocks$^3$, $m=20$, $B=4$, this gives $q=189$ queries for $128$-bit interactive soundness. The binding argument for the polynomial-valued chain is proximity-gap free: it uses only injectivity of the formal fold. Its price is a proof of size $\Theta(2^m)$. Specializing $X:=1$ collapses the protocol to classical sumcheck interleaved with scalar FRI; for this succinct variant we prove evaluation binding from an elementary correlated-agreement lemma for lines, valid for every linear code, giving the bound $m/|\K|+3n/|\K|+(3/4)^q$ at $B=4$, hence $q=309$. The two variants are validated numerically by an accompanying Rust implementation.
\end{abstract}

\tableofcontents

%======================================================================
\section{Setting and notation}
%======================================================================

\paragraph{Fields.} $k=\F_p$ with $p=2^{64}-2^{32}+1$ (Goldilocks), and $\K=\F_{p^3}$ a fixed cubic extension. All challenges are sampled uniformly from $\K$.

\paragraph{Parameters.} $m\ge1$ is the number of variables, $L=2^m$, $B\ge2$ is the blowup (a power of two), $n=BL$. We require $n\mid p-1$. The rate is $\rho=1/B$ and the unique-decoding threshold is
\[
\theta:=\frac{1-\rho}{2}.
\]

\paragraph{Domains.} For $0\le i\le m$ let $D_i=\mu_{B2^i}\subset k^\times$ be the group of $B2^i$-th roots of unity, so $|D_i|=B2^i$ and $x\mapsto x^2$ maps $D_i$ two-to-one onto $D_{i-1}$. The shrinking domains are $E_j:=D_{m-j}$. The fibers of the top domain over $D_0$ are
\[
\Phi_z:=\{y\in D_m:\ y^L=z\},\qquad z\in D_0,\qquad |\Phi_z|=L .
\]

\paragraph{Index conventions.} For $S\subseteq[m]$ put
\[
j(S):=\sum_{k\in S}2^{k-1},\qquad j'(S):=\sum_{k\in S}2^{m-k}.
\]
Both $S\mapsto j(S)$ and $S\mapsto j'(S)$ are bijections $2^{[m]}\to\{0,\dots,L-1\}$. For an integer $0\le t<L$, $t_k$ denotes the bit of weight $2^{k-1}$, so $t=\sum_k t_k2^{k-1}$. For $1\le i\le m$ and $0\le t<2^i$ let $\rev_i(t):=\sum_{k=1}^{i}t_k2^{i-k}$.

\paragraph{Polynomials.} The multilinear $f=\sum_{S\subseteq[m]}c_S\,X^S$ (with $X^S=\prod_{k\in S}X_k$) has Kronecker encoding
\[
F(Y):=f(Y,Y^2,\dots,Y^{2^{m-1}})=\sum_{S}c_S\,Y^{j(S)},\qquad \deg F<L.
\]
This is a $\K$-linear bijection between multilinear polynomials and $\K[Y]_{<L}$; we write $f=\kappa^{-1}(F)$. For $r=(r_1,\dots,r_m)\in\K^m$ define
\[
\beta_k(X):=r_kX^{2^{m-k}},\qquad f(\beta(X)):=f(\beta_1(X),\dots,\beta_m(X))=\sum_S c_S\Big(\prod_{k\in S}r_k\Big)X^{j'(S)}.
\]
For $0\le i\le m$ let $A_i(X_1,\dots,X_i):=\sum_{b\in\{0,1\}^{m-i}}f(X_1,\dots,X_i,b)$, so $A_m=f$ and $A_0=\sum_bf(b)$.

\paragraph{Even/odd split.} Every $P\in\K[X]$ is uniquely $P(X)=E_P(X^2)+X\,O_P(X^2)$.

\paragraph{Codes.} $\RS:=\RS[\K,D_m,L]=\{F|_{D_m}:\ F\in\K[Y]_{<L}\}$. Relative Hamming distance is $\dist(\cdot,\cdot)$. The minimum relative distance of $\RS$ is $\lambda=1-\rho+1/n>2\theta$.

\begin{remark}[The target is an evaluation]\label{rem:half}
Since $\sum_{b}\prod_{k\in S}b_k=2^{m-|S|}$,
\[
\sum_{b\in\{0,1\}^m}f(b)=\sum_S c_S\,2^{m-|S|}=2^m f\big(\tfrac12,\dots,\tfrac12\big).
\]
A hypercube sum of a single committed multilinear polynomial is therefore an evaluation claim at the point $(\frac12,\dots,\frac12)$. This identity explains the structure found below: both chains are re-encodings of the coefficient vector $(c_S)_S$, and the sum is a fixed linear functional of it.
\end{remark}

%======================================================================
\section{The protocol}
%======================================================================

\subsection{Chain A: growing-domain Kronecker--Sumcheck}

Oracles: $P_0:=H$ (a scalar); for $1\le i\le m$, tables $U_i,V_i:D_{i-1}\to\K$ and $P_i:D_i\to\K$.

Round $i$: the prover commits to $(U_i,V_i)$; the verifier samples $r_i$; the prover commits to $P_i$. The intended relations are
\[
P_{i-1}(Y)=2U_i(Y)+V_i(Y),\qquad P_i(X)=U_i(X^2)+r_iX\,V_i(X^2).
\]
\textbf{Path check.} For a point $x_m\in D_m$ let $x_{i-1}:=x_i^2$. The verifier opens, for $i=1,\dots,m$:
\begin{align*}
&(\mathrm a_i)\quad P_{i-1}(x_{i-1})\overset?=2U_i(x_{i-1})+V_i(x_{i-1})\qquad(P_0:=H),\\
&(\mathrm b_i)\quad P_i(x_i)\overset?=U_i(x_{i-1})+r_ix_i\,V_i(x_{i-1}).
\end{align*}

\subsection{Chain B: shrinking-domain polynomial-valued FRI}

\paragraph{Leaf format.} A level-$i$ leaf ($0\le i\le m$) sits at $z\in E_i=D_{m-i}$ and is an element
\[
s_i(z,\cdot)\in\K[W]_{<2^i},\qquad\text{representing } S_i(z,X):=s_i(z,X^{2^{m-i}}).
\]
It is committed as its canonical vector of $2^i$ coefficients. Level $0$ is the committed table $s_0(y)\in\K$, $y\in D_m$ (the oracle $C_f$).

\paragraph{Fold.} For $1\le i\le m$, $y\in E_{i-1}$, and $a,b\in\K[W]_{<2^{i-1}}$ put $e:=\frac{a+b}2$, $o:=\frac{a-b}{2y}$ and
\[
\Fold_{i,y}(a,b)(V):=e(V^2)+r_iV\,o(V^2)\ \in\K[V]_{<2^i}.
\]
With $V=X^{2^{m-i}}$ and $W=V^2=X^{2^{m-i+1}}$, this is exactly
\[
S_i(y^2,X)=\frac{S_{i-1}(y,X)+S_{i-1}(-y,X)}2+r_iX^{2^{m-i}}\,\frac{S_{i-1}(y,X)-S_{i-1}(-y,X)}{2y}.
\]
In round $i$, after $r_i$, the prover commits to the level-$i$ table $z\mapsto s_i(z,\cdot)$ on $E_i$.

\begin{lemma}[Injectivity of the fold]\label{lem:inj}
If $r_i\ne0$ and $y\ne0$, then $\Fold_{i,y}:\K[W]_{<2^{i-1}}^2\to\K[V]_{<2^i}$ is a linear bijection.
\end{lemma}
\begin{proof}
$e(V^2)$ is supported on even powers and $r_iVo(V^2)$ on odd powers, so the image determines $(e,o)$; then $a=e+yo$, $b=e-yo$. Dimensions agree ($2\cdot2^{i-1}=2^i$).
\end{proof}

\paragraph{Chain-B query.} For $y_0\in D_m$ let $y_{i}:=y_{i-1}^2\in E_i$. The verifier opens $s_{i}(\pm y_{i})$ and $s_{i+1}(y_{i+1})$ for $0\le i<m$ (pairs $\{y,-y\}$ are grouped in one Merkle leaf), checks the leaf format, and checks $s_{i+1}(y_{i+1})=\Fold_{i+1,y_i}\big(s_i(y_i),s_i(-y_i)\big)$.

\subsection{The repaired protocol}

The three items marked (F1)--(F3) are not optional; Section~\ref{sec:attacks} shows what fails without (F1) and why (F2) is load-bearing.

\begin{protocol}[$\Pi_{\mathrm{dual}}$]\label{prot:dual}\
\begin{enumerate}[label=\arabic*.,leftmargin=*]
\item Prover sends $C_f$ (oracle $s_0:D_m\to\K$) and $H$.
\item For $i=1,\dots,m$: prover commits $(U_i,V_i)$ on $D_{i-1}$; verifier sends $r_i\leftarrow\K$; prover commits $P_i$ on $D_i$ and the chain-B level $s_i$ on $E_i$.
\item Query phase (all commitments fixed):
\begin{enumerate}[label=(\alph*)]
\item[(F1)] Open all $B$ terminal leaves $s_m(z,\cdot)$, $z\in D_0$, in full; reject unless they are equal. Call the common value $T\in\K[X]_{<L}$.
\item[(F2)] Every opened leaf must be a coefficient vector of length exactly $2^i$ at level $i$.
\item[(F3)] Run $q$ chain-B queries from independent uniform $y_0^{(1)},\dots,y_0^{(q)}\in D_m$.
\item[(A)] Run one chain-A path check from a point $x_m\in D_m$ (any choice, e.g.\ $x_m=y_0^{(1)}$), and check $P_m(x_m)=T(x_m)$.
\end{enumerate}
\item Accept iff all checks pass.
\end{enumerate}
\end{protocol}

%======================================================================
\section{Completeness}
%======================================================================

\begin{lemma}[Honest chain A]\label{lem:chainA}
Define, for $1\le i\le m$,
\[
Q_i(Y,Z):=A_i\big(r_1Y^{2^{i-2}},r_2Y^{2^{i-3}},\dots,r_{i-1}Y,\ Z\big),\qquad Q_1(Y,Z):=A_1(Z),
\]
write $Q_i=U_i(Y)+Z\,V_i(Y)$, and set $P_i(X):=A_i(r_1X^{2^{i-1}},\dots,r_{i-1}X^2,r_iX)$, $P_0=A_0$. Then for all $i$:
\begin{enumerate}[label=(\roman*)]
\item $Q_i$ is affine in $Z$, and $\deg U_i,\deg V_i\le2^{i-1}-1$;
\item $P_{i-1}(Y)=2U_i(Y)+V_i(Y)$;
\item $P_i(X)=U_i(X^2)+r_iX\,V_i(X^2)$;
\item $P_m(X)=f(\beta(X))$.
\end{enumerate}
\end{lemma}
\begin{proof}
(i) $A_i$ is multilinear; the $Y$-degree is at most $\sum_{t=0}^{i-2}2^t=2^{i-1}-1$. (ii) $Q_i(Y,0)+Q_i(Y,1)=A_{i-1}(r_1Y^{2^{i-2}},\dots,r_{i-1}Y)=P_{i-1}(Y)$, using $A_{i-1}(\cdot)=A_i(\cdot,0)+A_i(\cdot,1)$. (iii) $Q_i(X^2,r_iX)=A_i(r_1X^{2^{i-1}},\dots,r_{i-1}X^2,r_iX)=P_i(X)$. (iv) is (iii)'s definition at $i=m$, since $A_m=f$.
\end{proof}

\begin{lemma}[Honest chain B]\label{lem:chainB}
Let $F=\sum_{t<L}c_tY^t$ and $s_0=F|_{D_m}$, and fold honestly. Then for $0\le i\le m$ and $z\in E_i$,
\[
s_i(z,V)=\sum_{t<L}c_t\ z^{\lfloor t/2^i\rfloor}\Big(\prod_{k\le i,\ t_k=1}r_k\Big)\,V^{\rev_i(t\bmod2^i)} .
\]
In particular $\deg_Vs_i<2^i$, and $S_m(z,X)=s_m(z,X)=f(\beta(X))$ for every $z\in D_0$.
\end{lemma}
\begin{proof}
Induction on $i$; $i=0$ is the definition. Assume the formula at level $i-1$ and fix $y\in E_{i-1}$, $z=y^2$. Write $u=\lfloor t/2^{i-1}\rfloor$, so $t_i=u\bmod2$. Then $\frac{y^u+(-y)^u}{2}=y^u[t_i=0]$ and $\frac{y^u-(-y)^u}{2y}=y^{u-1}[t_i=1]$, and in both cases $y^{u-t_i}=z^{\lfloor t/2^i\rfloor}$. The term with $t_i=0$ goes to $V^{2\rev_{i-1}(t)}$, the term with $t_i=1$ to $r_iV^{1+2\rev_{i-1}(t)}$, and $\rev_i(t\bmod2^i)=2\rev_{i-1}(t\bmod2^{i-1})+t_i$. At $i=m$: $\lfloor t/L\rfloor=0$ and $X^{\rev_m(j(S))}=X^{j'(S)}$.
\end{proof}

\begin{corollary}[Completeness]
The honest prover is accepted with probability $1$: all leaves are well formed, all $B$ terminal leaves equal $f(\beta)=P_m$, and every local check is an instance of Lemma~\ref{lem:chainA} or~\ref{lem:chainB}.
\end{corollary}

\begin{remark}[Orientation of the powers]
With the tempting definition
\[
Q_i(Y,Z)=A_i\big(r_1Y^{2^{i-1}},\dots,r_{i-1}Y^2,Z\big),
\]
identity (ii) would require the first of the following, whereas Lemma~\ref{lem:chainA} gives the second:
\[
P_{i-1}(Y)\overset?=A_{i-1}\big(r_1Y^{2^{i-1}},\dots,r_{i-1}Y^2\big),\qquad P_{i-1}(Y)=A_{i-1}\big(r_1Y^{2^{i-2}},\dots,r_{i-1}Y\big);
\]
the two differ for $i\ge2$ in general. Lemma~\ref{lem:chainA} is the consistent orientation.
\end{remark}

%======================================================================
\section{The exact-identity model}
%======================================================================

Suppose every round relation holds as a polynomial identity, all $r_i\neq0$, and the leaves have the format (F2).

\begin{proposition}[Chain A is invertible]\label{prop:Lambda}
The map $(U_i,V_i)\mapsto P_i$ with $P_{i-1}=2U_i+V_i$ is inverted by $U_i=E_{P_i}$, $V_i=r_i^{-1}O_{P_i}$, so $P_{i-1}=2E_{P_i}+r_i^{-1}O_{P_i}$. Iterating,
\[
H=P_0=\Lambda_r(P_m),\qquad
\Lambda_r\Big(\sum_{S}p_S\,X^{j'(S)}\Big):=\sum_S p_S\,2^{m-|S|}\prod_{k\in S}r_k^{-1}.
\]
If $P_m=f(\beta)$, i.e.\ $p_S=c_S\prod_{k\in S}r_k$, then $H=\sum_Sc_S2^{m-|S|}=\sum_bf(b)$.
\end{proposition}
\begin{proof}
At round $i$ the lowest power $X^1$ of $P_i$ carries variable $X_i$, so the odd part corresponds to $i\in S$; the inverse step multiplies coefficients with $i\notin S$ by $2$ and those with $i\in S$ by $r_i^{-1}$. Compose over $i=m,\dots,1$ and use Remark~\ref{rem:half}.
\end{proof}

\begin{proposition}[Chain B is a bijective re-encoding]\label{prop:Bbij}
For fixed $r$ with all $r_k\ne0$, $f\mapsto f(\beta)$ is a bijection from multilinear polynomials onto $\K[X]_{<L}$, and by Lemma~\ref{lem:inj} each fiber $\Phi_z$ of level~$0$ is mapped bijectively onto the terminal leaf $s_m(z,\cdot)$.
\end{proposition}

\begin{remark}\label{rem:norandom}
In the exact-identity model, neither chain uses the randomness of $r$: chain A computes the fixed functional $\Lambda_r$, which on honest inputs is the $r$-independent value $\sum_bf(b)$, and chain B is a deterministic butterfly of $F$. In particular a prover cannot satisfy the recurrences with an alternative polynomial family, since the recurrences are invertible. All soundness content therefore lies in the passage from sampled checks to identities, which is where the query analysis of Section~\ref{sec:sound} lives.
\end{remark}

%======================================================================
\section{Attacks on under-specified variants}\label{sec:attacks}
%======================================================================

\subsection{Fiber splitting: false sum accepted with probability one}

Let $\Pi^-$ be $\Pi_{\mathrm{dual}}$ with (F1) replaced by: \emph{every check that needs a terminal value at $x_m$ (resp.\ at the end of a chain-B query) uses the leaf $s_m(x_m^{L},\cdot)$ (resp.\ $s_m(y_m,\cdot)$) of its own fiber.}

Observe that both chains split into $B$ independent subtrees indexed by $z\in D_0$: in chain B, $\pm y$ and $y^2$ lie over the same $z$; in chain A, $x$ and $x^2$ satisfy $x^{2^i}=(x^2)^{2^{i-1}}$.

\begin{proposition}\label{prop:fiber}
For every multilinear $f_\star$ and every $H'\neq\sum_bf_\star(b)$ there is a prover for $\Pi^-$ whose initial table satisfies $\dist(s_0,F_\star)=B/n=1/L$ (so $F_\star=\kappa(f_\star)$ is its unique decoding) and which is accepted with probability $1$ claiming $H'$.
\end{proposition}
\begin{proof}
For $z\in D_0$ pick $y_z\in\Phi_z$ and set
\[
h_z(Y):=\frac{Y^L-z}{Y-y_z}=\sum_{t=0}^{L-1}y_z^{\,L-1-t}Y^t .
\]
It vanishes on $\Phi_z\setminus\{y_z\}$ and $h_z(y_z)=Ly_z^{L-1}\ne0$. Let $\eta_z:=\kappa^{-1}(h_z)$. By Remark~\ref{rem:half},
\[
\sigma_z:=\sum_b\eta_z(b)=\sum_{t<L}y_z^{L-1-t}2^{m-\wt(t)}=y_z^{L-1}\prod_{k=1}^m\big(2+y_z^{-2^{k-1}}\big).
\]
For each $k$ at most $2^{k-1}$ elements of the coset $\Phi_z$ satisfy $y^{2^{k-1}}=-\frac12$, so at most $L-1<|\Phi_z|$ choices are excluded, and we may choose $y_z$ with $\sigma_z\ne0$. Put $\lambda_z:=(H'-H_\star)/\sigma_z$ with $H_\star:=\sum_bf_\star(b)$, and
\[
s_0(y):=F_\star(y)+\lambda_zh_z(y)\qquad(y\in\Phi_z).
\]
This differs from $F_\star$ exactly at the $B$ points $y_z$. On subtree $z$ the prover runs the honest prover (both chains, Lemmas~\ref{lem:chainA} and~\ref{lem:chainB}) for $g_z:=f_\star+\lambda_z\eta_z$, whose Kronecker encoding interpolates $s_0$ on $\Phi_z$. Every local check holds identically, the leaf $s_m(z,\cdot)=g_z(\beta)$ matches chain A on the same subtree, and $\sum_bg_z(b)=H_\star+\lambda_z\sigma_z=H'$ for every $z$, consistent with the single scalar $P_0=H'$.
\end{proof}

Under (F1) this prover is rejected with certainty: $g_z(\beta)=g_{z'}(\beta)$ would force $\lambda_zh_z=\lambda_{z'}h_{z'}$, but $h_z$ vanishes on $\Phi_z\setminus\{y_z\}$ while $h_{z'}(y)=(z-z')/(y-y_{z'})\ne0$ there.

\subsection{Late commitment of the first round message}

\begin{proposition}\label{prop:late}
If $(U_1,V_1)$ is committed after $r_1$, then for any $H'\ne H_\star$ there is a prover that is honest in chain B and passes each independent uniform chain-A path with probability $1/2$.
\end{proposition}
\begin{proof}
For each $y\in D_0$ choose a root $x_y\in D_1$ with $2r_1x_y\ne1$ and solve $2U+V=H'$, $U+r_1x_yV=P_1^\star(x_y)$ (determinant $2r_1x_y-1$). Commit $P_1:=P_1^\star$ and play honestly from round $2$ on. Only $(\mathrm b_1)$ can fail, and only when $x_1=-x_y$; for uniform $x_m$, $x_1=x_m^{2^{m-1}}$ is uniform on $D_1$.
\end{proof}
With the correct order the corresponding error is at most $2B(2^m-1)/|\K|$, about $2^{-169}$ (Lemma~\ref{lem:pathA}), so the late variant drops chain A from field-size soundness to $2^{-q_A}$ for $q_A$ paths.

\subsection{Role of (F2) in the proof}
The detection argument (Lemma~\ref{lem:treeB}) walks up a query path and uses injectivity of $\Fold_{k,y}$ on the opened pair. On a query path only the path node is checked against its own children; the sibling is not. Lemma~\ref{lem:inj} gives injectivity only on $\K[W]_{<2^{k-1}}^2$, and without the format check a sibling could be an arbitrary polynomial of larger degree, for which $(a,b)\mapsto e(V^2)+r_iVo(V^2)$ has a nontrivial kernel. (F2) costs nothing and also fixes the leaf size at $2^i$ coefficients, which is what makes the cost bounds of Section~\ref{sec:cost} hold.

%======================================================================
\section{Soundness}\label{sec:sound}
%======================================================================

\subsection{The binding game}

\begin{definition}[Reference polynomial]
Given the committed oracle $s_0:D_m\to\K$, let $F_\star$ be the unique $F\in\K[Y]_{<L}$ with $\dist(s_0,F|_{D_m})<\theta$, if it exists (unique because $2\theta<\lambda$), and $f_\star:=\kappa^{-1}(F_\star)$, $H_\star:=\sum_bf_\star(b)$.
\end{definition}

\begin{definition}[Game]
An unbounded interactive oracle prover $\widetilde P$ interacts with the verifier of Protocol~\ref{prot:dual}. $\widetilde P$ wins if the verifier accepts and either $F_\star$ does not exist or $H\neq H_\star$.
\end{definition}

Note that $F_\star$ is a function of $s_0$ alone, hence fixed before $r_1$ is sampled.

\subsection{Chain A: pathwise sumcheck}

Let $U_i^\star,V_i^\star,P_i^\star$ be the honest tables of Lemma~\ref{lem:chainA} for $f_\star$; they depend on $f_\star$ and $r_1,\dots,r_{i-1}$ only (and $P_i^\star$ also on $r_i$).

\begin{lemma}[Pathwise sumcheck]\label{lem:pathA}
Assume $F_\star$ exists. Let $\Bad_A$ be the event that for some $i$, some $y\in D_{i-1}$ and some $x$ with $x^2=y$,
\[
(U_i(y),V_i(y))\neq(U_i^\star(y),V_i^\star(y))\quad\text{and}\quad U_i(y)+r_ixV_i(y)=U_i^\star(y)+r_ixV_i^\star(y).
\]
Then $\Pr[\Bad_A]\le\sum_{i=1}^m2|D_{i-1}|/|\K|=2B(2^m-1)/|\K|$. Moreover, outside $\Bad_A$, for \emph{every} $x_m\in D_m$: if $H\ne H_\star$ and all checks $(\mathrm a_i),(\mathrm b_i)$ on the path from $x_m$ pass, then $P_m(x_m)\neq f_\star(\beta(x_m))$.
\end{lemma}
\begin{proof}
Conditioned on the transcript before $r_i$, the tables $U_i,V_i,U_i^\star,V_i^\star$ are fixed and $r_i$ is uniform. For fixed $(y,x)$ with $(\Delta U,\Delta V):=(U_i-U_i^\star,V_i-V_i^\star)(y)\ne0$ and $x\ne0$, the affine function $r\mapsto\Delta U+rx\Delta V$ has at most one root. A union bound over $i$, $y$ and the two roots $x$ gives the probability bound.

For the second claim let $\Delta_i:=P_i(x_i)-P_i^\star(x_i)$, with $\Delta_0=H-H_\star\ne0$. If $\Delta_{i-1}\neq0$ and $(\mathrm a_i)$ passes, then $2\Delta U+\Delta V=\Delta_{i-1}\ne0$ at $x_{i-1}$, so $(\Delta U,\Delta V)\ne0$; outside $\Bad_A$ the right-hand side of $(\mathrm b_i)$ differs from $P_i^\star(x_i)$, and since $(\mathrm b_i)$ passes, $\Delta_i\ne0$. By induction $\Delta_m\ne0$, and $P_m^\star=f_\star(\beta)$ by Lemma~\ref{lem:chainA}(iv).
\end{proof}

\begin{remark}
Lemma~\ref{lem:pathA} needs no degree test on $P_i,U_i,V_i$, and one path suffices for any choice of $x_m$, including a fixed one. Wrongness propagates under pullback: a wrong value at $y$ forces wrong values at both square roots, so the error pattern of a false claim covers all of $D_m$, the opposite of FRI's distance decay.
\end{remark}

\subsection{Chain B: proximity-gap-free tree detection}

\begin{lemma}[Tree detection]\label{lem:treeB}
Assume all $r_k\ne0$, (F1) and (F2) pass with common terminal value $T$, and let $F'\in\K[Y]_{<L}$ be the unique polynomial with $\kappa^{-1}(F')(\beta)=T$ (Proposition~\ref{prop:Bbij}). Then every chain-B query started at $y_0$ with $s_0(y_0)\ne F'(y_0)$ rejects. Consequently
\[
\Pr\big[\text{all $q$ chain-B queries accept}\ \big|\ \text{commitments}\big]\le\big(1-\dist(s_0,F')\big)^q .
\]
\end{lemma}
\begin{proof}
Let $s_i'$ be the honest butterfly of $F'$ (Lemma~\ref{lem:chainB}); then $s_m'(z)=T$ for all $z\in D_0$, so $W_m=\emptyset$ where $W_i:=\{z\in E_i: s_i(z)\ne s_i'(z)\}$. Let $y_0\in W_0$ with ancestors $y_1,\dots,y_m$, and let $k\ge1$ be minimal with $y_k\notin W_k$. If the check at $y_k$ passed, then
\[
\Fold_{k,y_{k-1}}\big(s_{k-1}(y_{k-1}),s_{k-1}(-y_{k-1})\big)=s_k(y_k)=s'_k(y_k)=\Fold_{k,y_{k-1}}\big(s'_{k-1}(y_{k-1}),s'_{k-1}(-y_{k-1})\big),
\]
and since the opened leaves have the format (F2), Lemma~\ref{lem:inj} gives $s_{k-1}(y_{k-1})=s'_{k-1}(y_{k-1})$, contradicting $y_{k-1}\in W_{k-1}$. The queries are uniform and independent after $T$, hence $F'$, is fixed.
\end{proof}

\begin{lemma}[Distance to a wrong target]\label{lem:dist}
If $F_\star$ does not exist, or exists and $F'\ne F_\star$, then $\dist(s_0,F')\ge\theta$.
\end{lemma}
\begin{proof}
If $F_\star$ does not exist, $\dist(s_0,c)\ge\theta$ for every codeword $c$. Otherwise $\dist(s_0,F')\ge\lambda-\dist(s_0,F_\star)>\lambda-\theta>\theta$.
\end{proof}

\subsection{Main theorem}

\begin{theorem}[Interactive soundness of $\Pi_{\mathrm{dual}}$]\label{thm:main}
For every unbounded interactive oracle prover,
\[
\Pr[\Win]\ \le\ \frac{m}{|\K|}+\frac{2B(2^m-1)}{|\K|}+(1-\theta)^q .
\]
\end{theorem}
\begin{proof}
Partition the winning event.
\begin{enumerate}[label=(\arabic*)]
\item Some $r_k=0$: probability at most $m/|\K|$.
\item All $r_k\ne0$, and $F_\star$ does not exist or $T\ne f_\star(\beta)$. Then $F'$ is well defined and either $F_\star$ is absent or $F'\ne F_\star$, so by Lemmas~\ref{lem:treeB} and~\ref{lem:dist} acceptance has probability at most $(1-\theta)^q$.
\item All $r_k\ne0$, $F_\star$ exists, $T=f_\star(\beta)$, and $H\ne H_\star$. The check $P_m(x_m)=T(x_m)=f_\star(\beta(x_m))$ together with the chain-A path checks contradicts Lemma~\ref{lem:pathA} unless $\Bad_A$ holds.\qedhere
\end{enumerate}
\end{proof}

\subsection{Parameters}

For $\K=\F_{p^3}$ ($|\K|\approx2^{192}$), $m=20$, $B=4$ ($\theta=3/8$):

\begin{center}
\begin{tabular}{@{}lll@{}}
\toprule
Term & Value & $\log_2$ \\
\midrule
$m/|\K|$ & $20/p^3$ & $-187.7$ \\
$2B(2^m-1)/|\K|$ & $8(2^{20}-1)/p^3$ & $-169.0$ \\
$(5/8)^{189}$ & & $-128.16$ \\
$(5/8)^{188}$ & & $-127.48$ \\
\bottomrule
\end{tabular}
\end{center}

Hence $q=189$ chain-B queries and a single chain-A path give total interactive error below $2^{-128}$, and $q=189$ is minimal for this bound. This coincides with the unique-decoding query count of the line-binding symbolic-FRI protocol, as expected: both are governed by $(1-\theta)^q$.

\begin{remark}[Compilation]
Under the BCS transform with Merkle commitments in the random-oracle model, add $\varepsilon_{\mathrm{Merkle}}$, the probability that the prover finds a collision (at most $Q^2/2^{\lambda_h}$ for $Q$ queries and $\lambda_h$-bit digests). For the Fiat--Shamir NIZK, the query phase is a single round with error $(1-\theta)^q$ per attempt, so a prover with $Q$ hash queries gains a factor of about $Q$; compensate with grinding or with $q\ge(128+\log_2Q-g)/\log_2(8/5)$ for $g$ bits of proof-of-work.
\end{remark}

\begin{remark}[Beyond unique decoding]
Lemma~\ref{lem:treeB} gives per-query detection $\dist(s_0,F')$ for the target $F'$ chosen by the prover. Above $\theta$ the target may be any codeword of the list around $s_0$, selected adaptively after $r$, so the winning condition must change to a list statement; the deterministic chain alone cannot bind to a single $f_\star$ there.
\end{remark}

%======================================================================
\section{Costs}\label{sec:cost}
%======================================================================

All counts are in elements of $\K$ ($24$ bytes); $m=20$, $B=4$, $q=189$.

\paragraph{Prover.}
\begin{itemize}[leftmargin=*]
\item Chain B: level $i$ has $B2^{m-i}$ leaves of $2^i$ coefficients, i.e.\ $BL$ elements per level, $mBL\approx8.39\cdot10^7$ elements ($\approx2.0$ GB) in total, computed in $O(mBL)$ field operations.
\item Chain A: $\sum_{i=1}^m2B2^i\approx4BL\approx1.68\cdot10^7$ elements; the coefficients of all $A_i$ come from $O(L)$ partial sums, then one NTT per table.
\end{itemize}

\paragraph{Proof.} At level $i<m$ the opened data is at most $\min(q,B2^{m-i-1})\cdot2^{i+1}$ elements.
\begin{center}
\begin{tabular}{@{}lr@{}}
\toprule
Levels & Opened elements (upper bound) \\
\midrule
$0$--$13$ & $189\,(2^{15}-2)\approx6.19\cdot10^6$ \\
$14$--$19$ (whole levels) & $6\cdot BL\approx2.52\cdot10^7$ \\
terminal leaves (F1) & $BL\approx4.19\cdot10^6$ \\
chain A (one path) & $3m=60$ \\
\midrule
total & $\approx3.56\cdot10^7$ ($\approx0.85$ GB) \\
\bottomrule
\end{tabular}
\end{center}

\paragraph{Verifier.} $\Theta(L)$: reading $T$, evaluating $T(x_m)$, and fold checks whose cost equals the opened data. Once $T$ is read, the verifier could compute $H=\Lambda_r(T)$ directly, so in $\Pi_{\mathrm{dual}}$ chain A is redundant.

%======================================================================
\section{Collapse: specializing the formal variable}
%======================================================================

\begin{lemma}[Specialization]\label{lem:spec}
Fix $x\in k^\times$ and put $\alpha_i:=r_ix^{2^{m-i}}$. Then:
\begin{enumerate}[label=(\roman*)]
\item the scalars $g_i(z):=S_i(z,x)$ satisfy the scalar FRI fold
\[
g_i(y^2)=\frac{g_{i-1}(y)+g_{i-1}(-y)}2+\alpha_i\,\frac{g_{i-1}(y)-g_{i-1}(-y)}{2y};
\]
\item along the chain-A path from $x_m=x$, with $(u_i,v_i):=(U_i,V_i)(x_{i-1})$, the checks $(\mathrm a_i),(\mathrm b_i)$ read
\[
P_{i-1}(x_{i-1})=2u_i+v_i,\qquad P_i(x_i)=u_i+\alpha_iv_i,
\]
which is classical sumcheck with round polynomial $u_i+Zv_i$ and challenge $\alpha_i$;
\item the $\alpha_i$ are independent and uniform in $\K$.
\end{enumerate}
\end{lemma}
\begin{proof}
(i) and (ii) are substitution of $X=x$; note $x_i=x^{2^{m-i}}$ so $r_ix_i=\alpha_i$. (iii) $r_i\mapsto r_ix^{2^{m-i}}$ is a bijection of $\K$.
\end{proof}

The weighting $X^{2^{m-i}}$ is therefore cosmetic once $X$ is specialized, and the same $\alpha_i$ drives both chains. Taking $x=1$:

\begin{protocol}[$\Pi_1$]\label{prot:one}
Commit $s_0=F|_{D_m}$ and send $H$. For $i=1,\dots,m$: send $(u_i,v_i)\in\K^2$ in the clear, check $2u_i+v_i=c_{i-1}$ ($c_0=H$), sample $r_i$, set $c_i=u_i+r_iv_i$, commit the scalar FRI layer $g_i$ on $D_{m-i}$ folded with $r_i$. Finally check that $g_m$ is constant on $D_0$ with value $c_m$, and run $q$ FRI queries.
\end{protocol}

$\Pi_1$ is classical sumcheck for $\sum_bf(b)$ interleaved with scalar FRI on the Kronecker encoding, with shared challenges. The scalar fold is not injective, so Lemma~\ref{lem:treeB} does not transfer; the binding argument below replaces injectivity by an elementary correlated-agreement lemma for lines. Everything in this section holds verbatim for any family of linear codes $C_i\subseteq\K^{E_i}$ with the even/odd decomposition of Lemma~\ref{lem:pair}; Reed--Solomon is used only through that lemma and the value of the minimum distance.

\subsection{Pair representation}

Throughout, $g_0:=s_0$ and $g_i:E_i\to\K$ ($1\le i\le m$) are the committed layers of $\Pi_1$, $n_i:=|E_i|=n/2^i$, $L_i:=L/2^i$, and $C_i:=\RS[\K,E_i,L_i]$, with minimum distance $d_i=n_i-L_i+1$ and relative distance $\lambda_i=d_i/n_i>1-\rho$.

\begin{definition}[Pairs]
For $g:E_{i-1}\to\K$ and $z=y^2\in E_i$ put
\[
e_g(z):=\frac{g(y)+g(-y)}2,\qquad o_g(z):=\frac{g(y)-g(-y)}{2y},\qquad \Fold_r(g):=e_g+r\,o_g:E_i\to\K .
\]
The map $(g(y),g(-y))\mapsto(e_g(z),o_g(z))$ is invertible. The \emph{pair distance} $\Delta(g,g')$ is the fraction of $z\in E_i$ with $(e_g,o_g)(z)\ne(e_{g'},o_{g'})(z)$, equivalently $(g(\pm y))\ne(g'(\pm y))$.
\end{definition}

\begin{lemma}[Even/odd decomposition]\label{lem:pair}
For $1\le i\le m$:
\begin{enumerate}[label=(\roman*)]
\item $G\in C_{i-1}$ iff $(e_G,o_G)\in C_i\times C_i$, and then $G(Y)=e_G(Y^2)+Y\,o_G(Y^2)$ and $\Fold_r(G)\in C_i$;
\item $\dist(g,g')\le\Delta(g,g')$;
\item $\dist(\Fold_r g,\Fold_r g')\le\Delta(g,g')$;
\item for $F\in\K[Y]_{<L}$, $\Fold_{r_m}\circ\cdots\circ\Fold_{r_1}(F|_{D_m})$ is the constant $\kappa^{-1}(F)(r_1,\dots,r_m)$.
\end{enumerate}
\end{lemma}
\begin{proof}
(i) $G=E_G(Y^2)+YO_G(Y^2)$ with $\deg E_G,\deg O_G<L_i$, and $e_G=E_G|_{E_i}$, $o_G=O_G|_{E_i}$. (ii) A pair disagreement accounts for at most two point disagreements: $\dist(g,g')\le2\Delta(g,g')n_i/n_{i-1}=\Delta(g,g')$. (iii) The folds agree wherever the pairs agree. (iv) Lemma~\ref{lem:chainB} at $X=1$ (Lemma~\ref{lem:spec}(i) with $x=1$).
\end{proof}

\subsection{An elementary line lemma}

\begin{lemma}[Correlated agreement for lines]\label{lem:line}
Let $C\subseteq\K^N$ be linear with minimum distance $d$, let $t\ge0$ be an integer with $3t<d$, and let $e,o\in\K^N$. Put
\[
R:=\{r\in\K:\ \exists c\in C,\ |\{x:\ e_x+ro_x\ne c_x\}|\le t\}.
\]
If $|R|\ge\max(3,t+2)$, then there are $E,O\in C$ with $|\{x:\ (e_x,o_x)\ne(E_x,O_x)\}|\le t$.
\end{lemma}
\begin{proof}
For $r\in R$ the codeword $c_r$ is unique since $2t<d$; let $A_r$ be its agreement set, $|A_r|\ge N-t$. Fix distinct $r_1,r_2\in R$ and set $O:=(c_{r_1}-c_{r_2})/(r_1-r_2)$, $E:=c_{r_1}-r_1O\in C$. On $A_{r_1}\cap A_{r_2}$ the two equations $e+r_jo=c_{r_j}$ give $o=O$, $e=E$. For any $r\in R$, the codewords $c_r$ and $E+rO$ agree on $A_{r_1}\cap A_{r_2}\cap A_r$, of size at least $N-3t>N-d$, hence $c_r=E+rO$.

Let $D:=\{x:(e_x,o_x)\ne(E_x,O_x)\}$. For $x\in D$ the affine function $r\mapsto(e_x-E_x)+r(o_x-O_x)$ is not identically zero, so it vanishes for at most one $r$; let $\kappa_r$ be the number of $x\in D$ where it vanishes at $r$, so $\sum_{r}\kappa_r\le|D|$. For $r\in R$, the disagreement set of $e+ro$ with $c_r=E+rO$ has size $|D|-\kappa_r\le t$. Summing over $R$: $|R|(|D|-t)\le|D|$. If $|D|\ge t+1$ this gives $|R|\le|D|/(|D|-t)\le t+1$, a contradiction.
\end{proof}

\begin{corollary}\label{cor:line}
If $e,o$ are fixed and $r$ is uniform in $\K$, and no $E,O\in C$ satisfy $|\{x:(e_x,o_x)\ne(E_x,O_x)\}|\le t$, then $\Pr_r[r\in R]\le\max(2,t+1)/|\K|$.
\end{corollary}

\subsection{Binding of scalar FRI}

Fix the radius
\[
\delta^\ast:=\frac{1-\rho}3,\qquad t_i:=\lceil\delta^\ast n_i\rceil-1 .
\]
Then ``$<\delta^\ast$'' and ``at most $t_i$ disagreements'' coincide on $E_i$, and $3t_i<(1-\rho)n_i<d_i$, so Lemma~\ref{lem:line} applies to every $C_i$; moreover $2\delta^\ast<\lambda_i$. For $B=4$: $\delta^\ast=\frac14$ and $t_i=L_i-1$.

\begin{definition}[Reference objects]\label{def:refs}
In a run of $\Pi_1$:
\begin{itemize}[leftmargin=*]
\item $F_\star$ is the unique $F\in\K[Y]_{<L}$ with $\dist(s_0,F|_{D_m})<\delta^\ast$, if it exists; $f_\star=\kappa^{-1}(F_\star)$.
\item $\varphi_0:=g_0$ and $\varphi_i:=\Fold_{r_i}(g_{i-1})$ for $1\le i\le m$ (the values the verifier recomputes).
\item $G_j$ ($0\le j<m$) is the unique $G\in C_j$ with $\Delta(g_j,G)<\delta^\ast$, if it exists (unique by Lemma~\ref{lem:pair}(ii) and $2\delta^\ast<\lambda_j$). It is a function of $g_j$, hence fixed before $r_{j+1}$.
\item For a codeword $H\in C_j$, $H^{\to m}:=\Fold_{r_m}\circ\cdots\circ\Fold_{r_{j+1}}(H)$, a constant.
\end{itemize}
\end{definition}

\begin{definition}[Bad rounds]
Round $j+1$ ($0\le j<m$) is \emph{bad} if either
\begin{enumerate}[label=(\alph*)]
\item $G_j$ does not exist but $\dist(\varphi_{j+1},C_{j+1})<\delta^\ast$; or
\item for some reference $H\in\{\Fold_{r_j}\circ\cdots\circ\Fold_{r_{i+1}}(G_i):\ 0\le i\le j,\ G_i\text{ exists}\}\subseteq C_j$ and some $z\in E_{j+1}$, the pairs of $g_j$ and $H$ differ at $z$ but $\Fold_{r_{j+1}}(g_j)(z)=\Fold_{r_{j+1}}(H)(z)$.
\end{enumerate}
\end{definition}

\begin{lemma}\label{lem:badrounds}
$\Pr[\text{some round is bad}]\le\dfrac{1}{|\K|}\sum_{j=1}^m(j+1)\,n_j\le\dfrac{3n}{|\K|}$.
\end{lemma}
\begin{proof}
All objects in the definition of round $j+1$ are fixed before $r_{j+1}$ is sampled. (a) has probability at most $\max(2,t_{j+1}+1)/|\K|\le n_{j+1}/|\K|$ by Corollary~\ref{cor:line} applied to $(e_{g_j},o_{g_j})$ and $C_{j+1}$. For (b), each of the at most $j+1$ references and each $z$ contribute at most one value of $r_{j+1}$, so at most $(j+1)n_{j+1}/|\K|$. Sum over $j$, and use $\sum_{i\ge1}(i+1)2^{-i}=3$.
\end{proof}

From now on fix a run with no bad round in which the final layer $g_m$ is the constant $c$ (otherwise the verifier rejects). For $0\le i\le m$ and a uniform $z\in E_i$ let $\mathrm{Pass}_i(z)$ be the event that the checks $g_k(w)=\varphi_k(w)$ hold at every node $w=z^{2^{k-i}}$, $i\le k\le m$. Then a query from uniform $y_0$ passes iff $\mathrm{Pass}_0(y_0)$ (at level $0$ the check $g_0=\varphi_0$ is vacuous).

\begin{lemma}[Propagation]\label{lem:prop}
Let $H\in C_i$ be $G_i$ or a fold of some $G_{i'}$, $i'\le i$ (or, for $i=m$, any constant), with $H^{\to m}=c$. Then for every $z\in E_i$ with $\varphi_i(z)\ne H(z)$, $\mathrm{Pass}_i(z)$ fails.
\end{lemma}
\begin{proof}
Downward induction on $i$. For $i=m$: $\varphi_m(z)\ne H(z)=c=g_m(z)$, so the level-$m$ check fails. For $i<m$: if $\varphi_i(z)\ne g_i(z)$ the level-$i$ check fails. Otherwise $g_i(z)\ne H(z)$, so the pairs of $g_i$ and $H$ differ at $w=z^2$; since round $i+1$ is not bad (clause (b), with $H$ among the references), $\varphi_{i+1}(w)\ne\Fold_{r_{i+1}}(H)(w)$. The codeword $\Fold_{r_{i+1}}(H)$ is again a fold of some $G_{i'}$ with $(\Fold_{r_{i+1}}H)^{\to m}=c$, and the induction hypothesis at level $i+1$ shows that $\mathrm{Pass}_{i+1}(w)$, hence $\mathrm{Pass}_i(z)$, fails.
\end{proof}

\begin{lemma}[Detection]\label{lem:detect}
For $0\le i\le m$ let $H_i$ be the unique $H\in C_i$ with $\dist(\varphi_i,H)<\delta^\ast$, if it exists. If $H_i$ does not exist, or $H_i^{\to m}\ne c$, then $\Pr_z[\mathrm{Pass}_i(z)]\le1-\delta^\ast$.
\end{lemma}
\begin{proof}
Downward induction on $i$.

\emph{$i=m$.} $C_m$ is the code of constants on $D_0$ and $\Pr_z[\mathrm{Pass}_m(z)]=\Pr_z[\varphi_m(z)=c]=1-\dist(\varphi_m,c)$. If $H_m$ does not exist, $\dist(\varphi_m,c)\ge\delta^\ast$. If $H_m\ne c$, then $\dist(\varphi_m,c)\ge\lambda_m-\dist(\varphi_m,H_m)>\lambda_m-\delta^\ast>\delta^\ast$.

\emph{$i<m$.} Note $\Pr_z[\mathrm{Pass}_i(z)]\le\Pr_w[\mathrm{Pass}_{i+1}(w)]$, since $w=z^2$ is uniform on $E_{i+1}$ and $\mathrm{Pass}_i(z)\subseteq\mathrm{Pass}_{i+1}(z^2)$.
\begin{enumerate}[label=\arabic*.,leftmargin=*]
\item \emph{$G_i$ does not exist.} Since round $i+1$ is not bad (clause (a)), $\dist(\varphi_{i+1},C_{i+1})\ge\delta^\ast$, so $H_{i+1}$ does not exist; by induction $\Pr[\mathrm{Pass}_{i+1}]\le1-\delta^\ast$.
\item \emph{$G_i$ exists.} By Lemma~\ref{lem:pair}(iii), $\dist(\varphi_{i+1},\Fold_{r_{i+1}}G_i)\le\Delta(g_i,G_i)<\delta^\ast$, so $H_{i+1}=\Fold_{r_{i+1}}G_i$ and $H_{i+1}^{\to m}=G_i^{\to m}$.
\begin{enumerate}[label=2\alph*.]
\item If $G_i^{\to m}\ne c$: by induction $\Pr[\mathrm{Pass}_{i+1}]\le1-\delta^\ast$.
\item If $G_i^{\to m}=c$: by Lemma~\ref{lem:prop} with $H=G_i$, $\mathrm{Pass}_i(z)$ fails whenever $\varphi_i(z)\ne G_i(z)$, so $\Pr_z[\mathrm{Pass}_i(z)]\le1-\dist(\varphi_i,G_i)$. The hypothesis of the lemma says $H_i$ is absent or $H_i^{\to m}\ne c=G_i^{\to m}$, so in either case $H_i\ne G_i$; then $\dist(\varphi_i,G_i)\ge\delta^\ast$, directly if $H_i$ is absent and by $\dist(\varphi_i,G_i)\ge\lambda_i-\dist(\varphi_i,H_i)>\lambda_i-\delta^\ast>\delta^\ast$ otherwise.\qedhere
\end{enumerate}
\end{enumerate}
\end{proof}

\begin{theorem}[Evaluation binding of scalar FRI]\label{thm:bind}
For every prover in the FRI part of $\Pi_1$, with $q$ independent uniform queries,
\[
\Pr\big[\text{accept}\ \wedge\ (F_\star\text{ absent}\ \vee\ c\ne f_\star(r_1,\dots,r_m))\big]\ \le\ \frac{3n}{|\K|}+(1-\delta^\ast)^q .
\]
\end{theorem}
\begin{proof}
Outside the event of Lemma~\ref{lem:badrounds}, apply Lemma~\ref{lem:detect} at $i=0$: here $\varphi_0=g_0=s_0$, so $H_0=F_\star|_{D_m}$ when $F_\star$ exists, and $H_0^{\to m}=f_\star(r)$ by Lemma~\ref{lem:pair}(iv). Under the event in brackets each query passes with probability at most $1-\delta^\ast$, independently given the commitments.
\end{proof}

\subsection{Soundness of \texorpdfstring{$\Pi_1$}{Pi\_1}}

\begin{definition}[Game for $\Pi_1$]
$\widetilde P$ wins if the verifier accepts and either $F_\star$ (radius $\delta^\ast$) does not exist or $H\ne\sum_bf_\star(b)$.
\end{definition}

\begin{theorem}[Soundness of $\Pi_1$]\label{thm:one}
For every unbounded interactive oracle prover,
\[
\Pr[\Win]\ \le\ \frac{m}{|\K|}+\frac{3n}{|\K|}+(1-\delta^\ast)^q,\qquad\delta^\ast=\frac{1-\rho}3 .
\]
\end{theorem}
\begin{proof}
Outside the event of Theorem~\ref{thm:bind}, $F_\star$ exists and $c_m=c=f_\star(r)$. The sumcheck messages $(u_i,v_i)$ are sent before $r_i$, and the honest values $P_i^\star(1)=A_i^\star(r_1,\dots,r_i)$ depend on $f_\star$, which is fixed before $r_1$; Lemma~\ref{lem:pathA} with the single path through $x=1$ (one pair $(y,x)=(1,1)$ per round) bounds the event $H\ne H_\star$ by $m/|\K|$.
\end{proof}

\paragraph{Parameters.} For $B=4$: $\delta^\ast=\frac14$, $\log_2\frac43=0.41504$, and $(3/4)^{309}=2^{-128.25}$ while $(3/4)^{308}=2^{-127.83}$. With $3n/|\K|\approx2^{-168.4}$ and $m/|\K|\approx2^{-187.7}$, $q=309$ queries give interactive error below $2^{-128}$.

\paragraph{Cost of $\Pi_1$.} $2m$ scalars of sumcheck, and per query one pair and one Merkle path per level. With pairs grouped in leaves and $32$-byte digests this is $\sum_{i=0}^{m-1}\big(48+32(m-i+1)\big)=8320$ bytes per query, about $2.57$ MB for $q=309$ before path deduplication.

\begin{remark}[Radius]
The radius $\delta^\ast=(1-\rho)/3$ comes from Lemma~\ref{lem:line}, which holds for every linear code. Reaching the unique-decoding radius $(1-\rho)/2$ for Reed--Solomon codes amounts to extending Lemma~\ref{lem:line} to $2t<d$, which needs the algebraic structure of RS codes (decoding the line $e+Zo$ over $\K(Z)$); with it, the argument above goes through unchanged and gives $q=189$, matching $\Pi_{\mathrm{dual}}$.
\end{remark}

%======================================================================
\section{Numerical validation}
%======================================================================

The crate \texttt{rust/} of the repository implements $\Pi_{\mathrm{dual}}$ over Goldilocks$^3$ with $\K=\F_p[x]/(x^3-2)$ (irreducible, since $2^{(p-1)/3}\ne1$), the provers of Section~\ref{sec:attacks}, and a chain-B tightness prover: $s_0$ equals $F'$ outside a negation-closed set of relative size $3/8$ where it equals an unrelated codeword, and the higher levels are the honest butterfly of $F'$. Then $\dist(s_0,\RS)=\theta$, so $F_\star$ does not exist and every acceptance is a win. The tests check Lemmas~\ref{lem:chainA}, \ref{lem:chainB}, \ref{lem:inj}, Proposition~\ref{prop:Lambda} and Remark~\ref{rem:half} for $m\le6$ as exact identities.

Measured acceptance at $m=5$, $B=4$, $n=128$, over $2000$ independent runs (fresh challenges and queries):
\begin{center}
\begin{tabular}{@{}lrr@{}}
\toprule
Prover & measured & predicted \\
\midrule
honest, (F1) on & $1.0000$ & $1$ \\
fiber splitting, (F1) off & $1.0000$ & $1$ \\
fiber splitting, (F1) on & $0.0000$ & $0$ \\
late $(U_1,V_1)$, $1$ chain-A path & $0.4920$ & $0.5$ \\
late $(U_1,V_1)$, $2$ paths & $0.2590$ & $0.25$ \\
late $(U_1,V_1)$, $4$ paths & $0.0720$ & $0.0625$ \\
tightness, $q=1$ & $0.6325$ & $0.6250$ \\
tightness, $q=2$ & $0.4020$ & $0.3906$ \\
tightness, $q=4$ & $0.1595$ & $0.1526$ \\
tightness, $q=8$ & $0.0260$ & $0.0233$ \\
\bottomrule
\end{tabular}
\end{center}
All deviations are within $1.8$ standard deviations $\sqrt{p(1-p)/2000}$. The tightness rows show that the bound $(1-\theta)^q$ of Lemma~\ref{lem:treeB} is attained.

%======================================================================
\section{Summary of the security relation}
%======================================================================

\begin{center}
\begin{tabular}{@{}p{3.6cm}p{5.3cm}p{5.3cm}@{}}
\toprule
 & $\Pi_{\mathrm{dual}}$ (formal $X$) & $\Pi_1$ ($X:=1$) \\
\midrule
Chain A & pathwise sumcheck, one path, $\le2^{-169}$ & classical sumcheck, $m/|\K|$ \\
Chain B binding & injective fold, deterministic, $(1-\theta)^q$, $\theta=\frac{1-\rho}2$ & line lemma, $3n/|\K|+(1-\delta^\ast)^q$, $\delta^\ast=\frac{1-\rho}3$ \\
Role of $r$ in chain B & none & essential \\
Terminal link & $B$ full leaves equal, then $P_m(x_m)=T(x_m)$ & final layer constant $=c_m$ \\
Queries ($128$ bits, $B=4$) & $189$ & $309$ \\
Proof size ($m=20$) & $\approx0.85$ GB & $\approx2.6$ MB \\
Verifier & $\Theta(L)$ & $O(q\,m)$ \\
\bottomrule
\end{tabular}
\end{center}

The dual-chain design separates into a part that needs no proximity theory but is not succinct (the formal butterfly, bound by injectivity) and a succinct part bound by an elementary line lemma (scalar FRI). Remark~\ref{rem:half} shows that, for a single multilinear polynomial, the sum is the evaluation $2^mf(\frac12,\dots,\frac12)$; the sumcheck layer becomes necessary once the summand is a product (for instance $f\cdot g$ or $f\cdot\mathrm{eq}$), which is the natural next target for this construction.

\begin{thebibliography}{9}
\bibitem{lfkn} C.~Lund, L.~Fortnow, H.~Karloff, N.~Nisan. Algebraic methods for interactive proof systems. \emph{J.\ ACM} 39(4), 1992.
\bibitem{fri} E.~Ben-Sasson, I.~Bentov, Y.~Horesh, M.~Riabzev. Fast Reed--Solomon interactive oracle proofs of proximity. ICALP 2018.
\bibitem{bcs} E.~Ben-Sasson, A.~Chiesa, N.~Spooner. Interactive oracle proofs. TCC 2016-B.
\end{thebibliography}

\end{document}
